%-------------------------
% Younghwan Chae - Resume (English)
% Based on jakegut/resume (MIT License)
%-------------------------

\documentclass[letterpaper,11pt]{article}

\usepackage{latexsym}
\usepackage[empty]{fullpage}
\usepackage{titlesec}
\usepackage{marvosym}
\usepackage[usenames,dvipsnames]{color}
\usepackage{verbatim}
\usepackage{enumitem}
\usepackage[hidelinks]{hyperref}
\usepackage{fancyhdr}
\usepackage[english]{babel}
\usepackage{tabularx}
\usepackage{geometry}
\usepackage{needspace}

\geometry{top=0.4in, bottom=0.4in, left=0.45in, right=0.45in}

\pagestyle{fancy}
\fancyhf{}
\fancyfoot{}
\renewcommand{\headrulewidth}{0pt}
\renewcommand{\footrulewidth}{0pt}

\urlstyle{same}
\raggedbottom
\raggedright
\setlength{\tabcolsep}{0in}

\titleformat{\section}{
  \vspace{-4pt}\scshape\raggedright\large
}{}{0em}{}[\color{black}\titlerule \vspace{-5pt}]

%-------------------------
% Custom commands
\newcommand{\resumeItem}[1]{\item\small{{#1 \vspace{-2pt}}}}

\newcommand{\resumeSubheading}[4]{
  \vspace{-2pt}\item
    \begin{tabular*}{0.97\textwidth}[t]{l@{\extracolsep{\fill}}r}
      \textbf{#1} & #2 \\
      \textit{\small#3} & \textit{\small #4} \\
    \end{tabular*}\vspace{-7pt}
}

\newcommand{\resumeSubSubheading}[2]{
    \item
    \begin{tabular*}{0.97\textwidth}{l@{\extracolsep{\fill}}r}
      \textit{\small#1} & \textit{\small #2} \\
    \end{tabular*}\vspace{-7pt}
}

\newcommand{\resumeProjectHeading}[2]{
    \item
    \begin{tabular*}{0.97\textwidth}{l@{\extracolsep{\fill}}r}
      \small#1 & #2 \\
    \end{tabular*}\vspace{-7pt}
}

\newcommand{\resumeSubItem}[1]{\resumeItem{#1}\vspace{-4pt}}

\renewcommand\labelitemii{$\vcenter{\hbox{\tiny$\bullet$}}$}

\newcommand{\resumeSubHeadingListStart}{\begin{itemize}[leftmargin=0.15in, label={}]}
\newcommand{\resumeSubHeadingListEnd}{\end{itemize}}
\newcommand{\resumeItemListStart}{\begin{itemize}[itemsep=0pt]}
\newcommand{\resumeItemListEnd}{\end{itemize}\vspace{-5pt}}

%-------------------------------------------
\begin{document}

%-----------HEADING-----------
\begin{center}
    \textbf{\Huge \scshape Younghwan Chae} \\ \vspace{3pt}
    \small
    \href{mailto:chyohw97@gmail.com}{\underline{chyohw97@gmail.com}} \quad$|$\quad
    \href{https://www.linkedin.com/in/younghwan-chae/}{\underline{linkedin.com/in/younghwan-chae}} \quad$|$\quad
    \href{https://github.com/younghwan91}{\underline{github.com/younghwan91}}
\end{center}

%-----------SUMMARY-----------
\section{Summary}
\small{
Physical AI and robotics ML engineer with a PhD in numerical optimization (Mechanical Engineering) and \textbf{11 years of R\&D} (4 in industry, 7 in graduate research and lecturing), hands-on with \textbf{cameras} and \textbf{FMCW 4D imaging radar}. At bitsensing, shipped \textbf{camera/radar fusion perception} to \textbf{mass production} (\textbf{$\sim$200 units in 8 countries}, fusion RMSE \textbf{$-$51\%}) and built 3D radar perception (\textbf{PV-RCNN}, \textbf{BEV} tracking, \textbf{TensorRT INT8} on NVIDIA Jetson) and deep-learning DoA super-resolution (\textbf{2$\times$} angular resolution). At Doosan Robotics, uses FEM-trained \textbf{NVIDIA PhysicsNeMo} surrogates (\textbf{$\sim$15,000$\times$} faster than FEM) as the objective in sanding-path optimization, solved scan-to-path generation as an optimization problem (2--4\,h $\to$ 5--15\,s), and contributed early to an \textbf{RF-DETR} depalletizing detector (LoRA fine-tuning, mAP@0.5 94.2\%). \textbf{10 patent applications}, \textbf{6 peer-reviewed publications}.
}

%-----------EXPERIENCE-----------
\section{Experience}
  \resumeSubHeadingListStart

    \resumeSubheading
      {Doosan Robotics}{}
      {Principal Data Scientist \& ML Engineer --- Physical AI Platform Team (newly formed)}{Aug 2025 -- Present}
      \resumeItemListStart
        \resumeItem{Solved robotic-sanding \textbf{coverage-path generation} from raw 3D scans as an optimization problem, replacing a 2--4\,h manual process with a 5--15\,s pipeline (\textbf{$>$99\% reduction}): RBF surface reconstruction (sub-0.5\,mm RMSE), \textbf{curvature-adaptive path optimization}, and surface-normal-aligned tool orientation (TNB frames); vectorized boundary masking (100--1,000$\times$ faster) and a 3D review dashboard with live RMSE/MAE/max error}
        \resumeItem{Integrated the path pipeline as an \textbf{agentic skill} in the team's \textbf{Agentic Robot OS} (rolling out), with the LLM agent calling the path optimizer as a tool. The surface fitting was developed in-house; in an \textbf{internal technical evaluation} it outperformed a competing external solution, contributing to a strategic sourcing decision}
        \resumeItem{Led the \textbf{physics-based surface-quality model} for robotic sanding and derived its equations single-handedly (contact pressure $\to$ accumulated work $\to$ roughness \textbf{Ra}) with constants calibrated from designed experiments; preprocessed data and generated \textbf{FEM} ground truth; also applied reinforcement learning to sanding path/motion optimization using the robot's URDF model}
        \resumeItem{\textbf{Surrogate-based path optimization}: optimized the \textbf{NVIDIA Newton} (VBD) simulator itself to reproduce the FEM ground truth (calibrating contact stiffness, damping, friction), used it to generate data at scale for \textbf{NVIDIA PhysicsNeMo} surrogates plus a residual surrogate on the FEM residual, cutting one pressure-field evaluation from $\sim$1\,h of FEM to $\sim$0.24\,s (\textbf{$\sim$15,000$\times$}). This work led to an \textbf{NVIDIA HQ collaboration workshop}, where I presented it}
        \resumeItem{Contributed to the early development of the \textbf{detection model} for an industrial depalletizing vision system handling \textbf{densely stacked, partially occluded} boxes: identified \textbf{RF-DETR} as a better fit than YOLO-family models for precision on overlapping boxes, trained it with \textbf{LoRA} fine-tuning and SAM segmentation, and handed it over to the owning team (validation set: \textbf{mAP@0.5 94.2\%}, mAP@0.5:0.95 78.6\%, precision 96.1\%, recall 91.8\%); built a dataset $\to$ training $\to$ evaluation-report pipeline}
        \resumeItem{Primary developer of a \textbf{single-RGB-camera imitation-learning data pipeline} that replaces teleoperation and motion-capture rigs: extracts human joints from demonstration video with \textbf{MediaPipe}, retargets them to the robot by solving \textbf{inverse kinematics at the TCP}, and outputs imitation-learning / \textbf{VLA} training data (data-collection dashboard and QA flow); collected demonstration data myself}
        \resumeItem{Took part in the \textbf{NVIDIA Cosmos Cookoff} using NVIDIA Cosmos and developed a pallet-pattern optimization demo for it: searched stacking patterns with \textbf{9 industrial pallet standards} (dimensions, max load) as constraints, optimizing area utilization, center-of-mass offset, and layer-height uniformity; 3D visualization and auto-generated local-LLM analysis reports (prompt tuning) and video; generated \textbf{MuJoCo} box-stack data (varied sizes and stacking combinations) and tested Cosmos augmentation (lighting, placement, object type)}
        \resumeItem{Built an \textbf{LLM + knowledge-graph} Q\&A system: pluggable ingestion of documents, source code, Jira, and Confluence into an OWL/RDF ontology; hybrid vector/graph/keyword retrieval with cited answers; logging of queries, retrievals, answers, and user feedback as evaluation data}
      \resumeItemListEnd

    \needspace{6\baselineskip}
    \resumeSubheading
      {bitsensing Inc.}{}
      {Technical Research Personnel (alternative military service) --- completed}{Aug 2022 -- Aug 2025}
    \resumeSubSubheading
      {AI Perception Engineer --- Signal Processing Team}{Sep 2023 -- Aug 2025}
      \resumeItemListStart
        \resumeItem{Developed \textbf{3D radar perception} for autonomous driving on 4D imaging-radar point clouds: detection (\textbf{PV-RCNN}), point-feature learning (\textbf{DGCNN}), and 3D MOT (\textbf{CenterPoint-BEV}, \textbf{AB3DMOT})}
        \resumeItem{\textbf{DoA estimation under uncertainty}: posed deep-learning super-resolution direction-of-arrival estimation on 4D imaging radar as an \textbf{inverse problem} (recovering the angular spectrum from undersampled measurements) and validated it in real-world experiments --- \textbf{2$\times$ angular resolution} over conventional FFT/beamforming}
        \resumeItem{Built ML pipelines for signal denoising, bias correction, and synthetic data generation; compressed models via \textbf{knowledge distillation}, pruning, and \textbf{TensorRT INT8 quantization} for real-time embedded deployment on \textbf{NVIDIA Jetson} (C++) without accuracy regression}
        \resumeItem{Filed \textbf{7 patent applications} in AI-based signal processing, uncertainty quantification, and embedded learning}
      \resumeItemListEnd

    \needspace{5\baselineskip}
    \resumeSubSubheading
      {Sensor Fusion Engineer --- Tracking and Application Team}{Aug 2022 -- Sep 2023}
      \resumeItemListStart
        \resumeItem{Built a camera/radar \textbf{fusion perception system}: \textbf{YOLO-based 8-class obstacle detection} (vehicle, pedestrian, motorcycle, etc.), lane \textbf{semantic segmentation} (PP-LiteSeg, U-Net), ML-based false-positive removal (radar FP/FN cases categorized), and \textbf{time synchronization} of radar/camera HDF5 recordings uploaded to the server for analysis and algorithm development}
        \resumeItem{Cut camera/radar fusion RMSE by \textbf{51\%} by solving the calibration as \textbf{non-linear reprojection-error minimization} with camera calibration and \textbf{multi-view geometry}; also trained multi-camera + radar BEV fusion models, both geometry-based and transformer-based (BEVFusion), estimating 3D boxes (position and heading), with LiDAR used only as ground truth}
        \resumeItem{\textbf{Large-scale multi-object tracking} of up to \textbf{256 targets across 12 lanes}: per-frame Hungarian-based SORT assignment combined with \textbf{IMM/Kalman} sequential state estimation, running on an embedded system; related paper on track-based radar/vision fusion (\textit{IEEE Sensors Journal} '25, co-first author)}
        \resumeItem{Shipped to \textbf{mass production} (\textbf{$\sim$200 radar units in 8 countries}); \textbf{TensorRT} edge-optimized; stable in adverse weather}
        \resumeItem{\textbf{Led} and productized auto-calibration motor control that holds the traffic radar's optimal field of view; used an IMU for sensor orientation and position}
        \resumeItem{Filed \textbf{3 patent applications} in sensor fusion and probabilistic object tracking}
      \resumeItemListEnd

    \needspace{5\baselineskip}
    \resumeSubheading
      {University of Pretoria}{Pretoria, South Africa}
      {Postgraduate Researcher \& Assistant Lecturer}{Jan 2015 -- Dec 2021}
      \resumeItemListStart
        \resumeItem{PhD: proposed \textbf{GOALS} (Gradient-Only Approximations for Line Searches), which sets step sizes from sign changes of the directional derivative alone, giving stable, consistent deep-network training on discontinuous, stochastic losses; published in \textit{Applied Intelligence} (Springer Nature)}
        \resumeItem{Developed \textbf{sub-dimensional surrogates} decomposing high-dimensional ML optimization under a limited evaluation budget into low-dimensional subspace approximations; published in \textit{Advances in AI Reviews} (IFSA)}
        \resumeItem{Formulated \textbf{optimal sensor placement} on FEM simulations as an inverse problem and a budget-constrained \textbf{combinatorial selection} over $\sim$$10^{26}$ candidate layouts, solved with a rank-variance heuristic; published in \textit{Engineering Structures} (Elsevier)}
        \resumeItem{BEng thesis (2014): \textbf{FEM-based topology optimization} of a compliant plier maximizing force-transmission efficiency; 3D-printed and validated with strain gauges}
        \resumeItem{\textbf{Assistant Lecturer}: taught Mathematical Optimization and Numerical Analysis in English to undergraduates}
      \resumeItemListEnd

  \resumeSubHeadingListEnd

%-----------EDUCATION-----------
\needspace{5\baselineskip}
\section{Education}
  \resumeSubHeadingListStart
    \item
    \begin{tabular*}{0.97\textwidth}[t]{l@{\extracolsep{\fill}}r}
      \textbf{University of Pretoria} & Pretoria, South Africa \\
      \textit{\small Doctor of Philosophy (PhD), Mechanical Engineering --- Numerical Optimization} & \textit{\small 2017 -- 2021} \\
      \textit{\small Master of Engineering (MEng), Mechanical Engineering} {\small --- Cum Laude} & \textit{\small 2016} \\
      \textit{\small Bachelor of Engineering Honours (BEng Hons), Mechanical Engineering} {\small --- Cum Laude} & \textit{\small 2015} \\
      \textit{\small Bachelor of Engineering (BEng), Mechanical Engineering} {\small --- Cum Laude} & \textit{\small 2011 -- 2014} \\
    \end{tabular*}\vspace{-7pt}
  \resumeSubHeadingListEnd

%-----------SKILLS-----------
\needspace{5\baselineskip}
\section{Technical Skills}
 \begin{itemize}[leftmargin=0.15in, label={}, itemsep=0pt, parsep=0pt, topsep=0pt]
    \small{\item{\textbf{Programming}{: Python, C/C++, MATLAB}}}
    \small{\item{\textbf{ML \& Scientific Computing}{: PyTorch, Transformers (RF-DETR), CNNs (YOLO), 3D/Point-Cloud (PV-RCNN, CenterPoint, DGCNN), SAM, LoRA, MediaPipe, NVIDIA PhysicsNeMo, NVIDIA Newton, Knowledge Distillation, TensorRT, TensorFlow, OpenCV, CUDA, ONNX, CVXPY, Stable-Baselines3, ML-Agents, MuJoCo, NumPy, SciPy, scikit-learn, pandas}}}
    \small{\item{\textbf{Sensors}{: FMCW Radar (4D-Imaging), Camera (RGB, PTZ), LiDAR (ground truth), IMU, GPS}}}
    \small{\item{\textbf{Domains}{: Numerical Optimization, FEM Simulation, Sensor Fusion, Camera/Radar Time Synchronization, Camera Calibration \& Multi-view Geometry, Signal Processing, 3D Perception, Point Cloud Processing, Semantic Segmentation, Self-supervised Learning, Probabilistic Modeling, Kalman Filter, Bayesian Inference, Surrogate Modeling, Reinforcement Learning}}}
    \small{\item{\textbf{Tools}{: Git, GitHub Actions CI, pytest/mypy/Playwright, uv/ruff, Docker, Linux/Unix, Jupyter, Dash/Plotly, DuckDB, MySQL, HDF5, Roboflow, Ollama, \LaTeX}}}
    \small{\item{\textbf{Languages}{: Korean (native), English (fluent)}}}
 \end{itemize}

%-----------PUBLICATIONS-----------
\needspace{5\baselineskip}
\section{Publications}
 \begin{itemize}[leftmargin=0.15in, label={}]
    \small{\item{
      \textbf{[IEEE Sensors Journal '25]}\enspace S. Kwak, \underline{Y. Chae}, M. Choi, et al. ``Improved accuracy of track-based integration of radar and vision sensors.'' \textit{IEEE Sensors Journal}, vol. 25, no. 10. \textit{(Co-first author)} \\[1pt]
      \textbf{[IEEE Sensors Journal '24]}\enspace C. Park, H.I. Baek, \underline{Y. Chae}, H.S. Lim, et al. ``Deep learning-based kick motion recognition in millimeter wave band radar system.'' \textit{IEEE Sensors Journal}, vol. 24, no. 19. \\[1pt]
      \textbf{[IEEE RadarConf '24]}\enspace H.I. Baek, \underline{Y. Chae}, H.S. Lim, et al. ``Implementation of deep learning-based kick gesture recognition using 60 GHz radar sensor.'' \\[1pt]
      \textbf{[Applied Intelligence '23]}\enspace \underline{Y. Chae}, D.N. Wilke, D. Kafka. ``Gradient-only surrogate to resolve learning rates for robust and consistent training of deep neural networks.'' \textit{Applied Intelligence} (Springer Nature). \\[1pt]
      \textbf{[Advances in AI Reviews '19]}\enspace \underline{Y. Chae}, D.N. Wilke. ``Sub-dimensional surrogates to solve high-dimensional optimization problems in machine learning.'' \textit{Advances in AI Reviews} (IFSA). \\[1pt]
      \textbf{[Engineering Structures '17]}\enspace \underline{Y. Chae}, D.N. Wilke. ``Heuristic linear algebraic rank-variance formulation for efficient sensor placement.'' \textit{Engineering Structures} (Elsevier). \\[6pt]
      \textsc{Preprints} \\[2pt]
      \textbf{[arXiv '21]}\enspace \underline{Y. Chae}, D.N. Wilke, D. Kafka. ``GOALS: Gradient-only approximations for line searches towards robust and consistent training of deep neural networks.'' arXiv:2105.10915. \\[1pt]
      \textbf{[arXiv '19]}\enspace \underline{Y. Chae}. ``Empirical study towards understanding line search approximations for training neural networks.'' arXiv:1909.06893.
    }}
 \end{itemize}

%-----------PATENTS-----------
\needspace{5\baselineskip}
\section{Patent Applications \normalfont{\small(10 filed, domestic \& international: US, EP)}}
 \begin{itemize}[leftmargin=0.15in, label={}]
    \small{\item{
      \textbf{2023}\enspace Sensor fusion for accurate environment perception; object recognition for unstructured task execution \\
      \textbf{2023}\enspace Sensor data structuring and utilization for efficient environment information processing \\
      \textbf{2023}\enspace Sensor data quality enhancement via AI/ML-based signal processing \\
      \textbf{2023}\enspace Data acquisition and model robustness enhancement for AI model training \\
      \textbf{2023}\enspace Sensor data processing and perception system performance improvement \textit{(US, EP)} \\
      \textbf{2024}\enspace Environment perception reinforcement and AI/ML-based sensor data analysis under adverse conditions \\
      \textbf{2024}\enspace Sensor data accuracy assurance and overall reliability of control systems \\
      \textbf{2025}\enspace Core sensor data processing for precise object localization and tracking \\
      \textbf{2025}\enspace High-quality training data acquisition for AI model learning \\
      \textbf{2025}\enspace Uncertainty quantification for reliable perception in uncertain environments
    }}
 \end{itemize}

%-----------AWARDS-----------
\needspace{5\baselineskip}
\section{Awards \& Honors}
 \begin{itemize}[leftmargin=0.15in, label={}]
    \small{\item{
      2025\enspace Top rating in the annual performance review (AI\&SW Division) --- Doosan Robotics \\
      2017\enspace Academic Honorary Colours, MEng --- University of Pretoria \\
      2016\enspace Merit Certificate for Academics, BEng Honours --- University of Pretoria \\
      2015\enspace Achievement Award \& Academic Honorary Colours, BEng --- University of Pretoria \\
      2011\enspace Golden Key International Honour Society \\
      2010\enspace South African National Mathematics Olympiad Top 100; 2nd in Province
    }}
 \end{itemize}

%-----------OPEN SOURCE & COMMUNITY-----------
\needspace{4\baselineskip}
\section{Open Source \& Community}
 \begin{itemize}[leftmargin=0.15in, label={}, itemsep=2pt]
    \small{\item{2026\enspace \textbf{API Ambassador}, Next Securities --- open-source quant pipelines on GitHub (data collection, backtesting, execution; \href{https://github.com/younghwan91}{\underline{github.com/younghwan91}}); developer interview ahead of the US-equity REST API beta, then an ambassador contract}}
 \end{itemize}

%-------------------------------------------
\end{document}
